\section{The direct algebraic-reduction proof}
\label{sec:algebraic-dimension}

The calculation in this section is independent of the integral topology of
\(X\).  It uses the period family of
\cref{prop:period-input-audit} and the proper connected fibration of
\cref{thm:candidate-global-geometry}; it does not use \(H^2(X;\mathbb Z)\),
the sphere-recognition argument, or any theorem of
Campana--Demailly--Peternell.

\begin{lemma}[algebraic reduction of a two-dimensional torus]
\label{lem:torus-algebraic-reduction-audit}
Let \(T\) be a complex torus of dimension two.  Suppose
\[
 \operatorname{NS}(T)=\mathbb Z\eta,\qquad \eta^2\ne0,
\]
and the Appell--Humbert Hermitian form \(H_\eta\) is nondegenerate and
indefinite.  Then \(a(T)=0\).
\end{lemma}

\begin{proof}
The algebraic reduction of a compact complex torus is a quotient
\(q:T\to A=T/K\) by a connected complex subtorus, with \(A\) an abelian
variety and \(\dim A=a(T)\)
\cite[Chapter~2, \S6]{BirkenhakeLange1999}; see also
\cite[\S8, p.~2437]{PareschiPopaSchnell2017}.

Assume \(a(T)>0\).  If \(a(T)=2\), then \(T\) is an abelian surface.  An
ample line bundle has first Chern class \(n\eta\), \(n\ne0\), and its
Appell--Humbert form \(nH_\eta\) is positive definite.  This contradicts
the indefiniteness of \(H_\eta\).

If \(a(T)=1\), then \(A\) is an elliptic curve.  For a point \(p\in A\),
put
\[
 \varphi=q^*c_1(\mathcal O_A(p))\in\operatorname{NS}(T).
\]
The class is nonzero: for any translation-invariant Kahler form \(\omega\)
on \(T\), the projection formula gives
\(\int_T\varphi\wedge\omega=\int_K\omega>0\).  But
\(\varphi^2=0\), since \(A\) is a curve.  Writing
\(\varphi=n\eta\), \(n\ne0\), yields
\(0=n^2\eta^2\ne0\), a contradiction.  Both possible positive algebraic
dimensions are excluded.
\end{proof}

Let
\[
 V=\mathbb Z\langle\gamma,u,w,\delta\rangle,
 \qquad
 \eta=u\wedge w+6\gamma\wedge\delta.
\]
For an integral two-class \(E\in\bigwedge^2V\), write its coordinates in
the ordered basis
\((\gamma u,\gamma w,\gamma\delta,uw,u\delta,w\delta)\) as
\(E_{\gamma u},\ldots,E_{w\delta}\), and set
\[
\begin{aligned}
P_E={}&E_{\gamma u}-E_{\gamma w}\tau-E_{\gamma\delta}\mu
 +6E_{uw}\mu+E_{u\delta}\beta\\
&+E_{w\delta}(6\mu^2-\tau\beta).
\end{aligned}
\tag{3.1}
\]

\begin{proposition}[very general Neron--Severi group]
\label{prop:very-general-torus-audit}
\Verified\ For the descended period family of
\cref{prop:period-input-audit}, outside a countable subset of
\(B^\circ\) the smooth fibre
\(F_b\) satisfies
\[
 \operatorname{NS}(F_b)=\mathbb Z\eta,
 \qquad a(F_b)=0.
\tag{3.2}
\]
\end{proposition}

\begin{proof}
The two holomorphic one-forms defining the Hodge subspace are
\[
 \sigma_1=6\mu\gamma+\tau u+w,
 \qquad
 \sigma_2=\beta\gamma+\mu u+\delta.
\]
The determinant inequality (2.2) gives
\(V_{\mathbb C}=F^1\oplus\overline{F^1}\), where
\(F^1=\langle\sigma_1,\sigma_2\rangle\).  For a real two-class \(E\),
the Hodge decomposition therefore gives
\[
 E\in H^{1,1}(F_b)
 \quad\Longleftrightarrow\quad
 E\wedge\sigma_1\wedge\sigma_2=0.
\]
Expanding the wedge product gives
\[
 E\wedge\sigma_1\wedge\sigma_2
 =P_E\,\gamma\wedge u\wedge w\wedge\delta.
\tag{3.3}
\]

We next determine exactly which \(P_E\) vanish identically.  Suppose
\[
 a+b\tau+c\mu+d\beta+e(6\mu^2-\tau\beta)\equiv0.
\tag{3.4}
\]
Apply the cusp transformation
\((\tau,\mu,\beta)\mapsto(\tau-1,\mu,\beta+1)\) and subtract.  This gives
\[
 -b+d+e(\beta-\tau+1)=0.
\]
Because \(\beta-\tau\) is translated by two, \(e=0\), and then \(d=b\).
Apply the order-four transformation in (2.1), subtract (3.4), and multiply
by \(-\tau\).  One obtains
\[
 6b\mu^2+c(\tau-1)\mu+b(\tau^2+3\tau+1)-c\tau=0.
\]
Subtract its cusp translate:
\[
 c\mu+2b(\tau+1)-c=0.
\tag{3.5}
\]
Translate (3.5) once around the cusp.  Since \(\mu\) is invariant and
\(\tau\mapsto\tau-1\), subtracting the translated identity from (3.5)
gives \(2b=0\).  Thus \(b=0\), and (3.5) becomes
\(c(\mu-1)=0\).  If \(c\ne0\), then \(\mu\equiv1\), whereas the
order-four law would give
\(1=\mu\circ g_2=1+1/\tau\), an impossibility.  Hence \(c=0\), then
\(b=d=0\), and finally \(a=0\).
Thus the five functions
\[
 1,\quad\tau,\quad\mu,\quad\beta,\quad6\mu^2-\tau\beta
\]
are linearly independent over \(\mathbb C\).

Reading the coefficients in (3.1) now gives
\[
 P_E\equiv0\quad\Longleftrightarrow\quad E\in\mathbb Z\eta.
\tag{3.6}
\]
For each other integral class, the zero set of \(P_E\) on the connected
Riemann surface \(\mathfrak h_z\) is discrete, hence countable.  There are
only countably many integral classes.  The class
\(\eta=q=uw+6\gamma\delta\) is monodromy invariant by
\cref{thm:coinvariants}, so it defines a class on every marked fibre after
descent.  The image in \(B^\circ\) of the countable union of the remaining
discrete zero sets is still countable.  A base point outside that image has
no additional integral \((1,1)\)-class on any lift, proving the
Neron--Severi statement on the quotient family.

Finally,
\[
 \eta^2=12\,\gamma\wedge u\wedge w\wedge\delta.
\tag{3.7}
\]
For a nondegenerate alternating form \(E\), write \(E^\vee\) for its
inverse alternating form on the dual real vector space.  If
\(G_\eta=(i\eta^\vee(\sigma_a,\overline{\sigma_b}))_{a,b}\), direct
inversion of the alternating-form matrix and substitution give
\[
G_\eta=
\begin{pmatrix}
2\operatorname{Im}\tau&2\operatorname{Im}\mu\\
2\operatorname{Im}\mu&\frac13\operatorname{Im}\beta
\end{pmatrix},
\]
\[
\det G_\eta=
\frac23\operatorname{Im}\tau
\left(\operatorname{Im}\beta-
\frac{6(\operatorname{Im}\mu)^2}{\operatorname{Im}\tau}\right)<0.
\tag{3.8}
\]
Here is the Appell--Humbert bridge needed to read the inertia of this
matrix.  With the convention \(E=\operatorname{Im}H_E\), one has
\[
 E=\frac1{2i}\sum_{a,b}(H_E)_{ab}\,
       \sigma_a\wedge\overline{\sigma_b}.
\]
Thus, in the ordered basis
\((\sigma_1,\sigma_2,\overline{\sigma_1},\overline{\sigma_2})\), the
upper-right block of the alternating-form matrix of \(E\) is
\(A=-(i/2)H_E\).  Inverting the block matrix gives
\[
 E^\vee(\sigma_a,\overline{\sigma_b})
   =[-(A^t)^{-1}]_{ab}
   =-2i(\overline{H_E}^{\,-1})_{ab},
 \qquad
 G_E=2\overline{H_E}^{\,-1}.
\tag{3.8a}
\]
The displayed \(G_\eta\) is real symmetric, so
\(H_\eta=2G_\eta^{-1}\).  Inversion and multiplication by a positive scalar
preserve inertia.  Hence \(H_\eta\) is nondegenerate and indefinite, while
(3.7) is nonzero.
Apply \cref{lem:torus-algebraic-reduction-audit}.
\end{proof}

\begin{lemma}[meromorphic descent]
\label{lem:meromorphic-descent-audit}
Let \(f:X\to B\) be a proper surjective holomorphic map from a connected
complex threefold to a connected compact Riemann surface, with connected
fibres, and a submersion over
\(B^\circ=B\setminus\Sigma\), where \(\Sigma\) is finite and
\(B^\circ\ne\varnothing\).  If
\(g\in\mathcal M(X)\) is constant on the fibres over an uncountable subset
of \(B^\circ\), then there is a unique \(\varphi\in\mathcal M(B)\) such that
\(g=f^*\varphi\).
\end{lemma}

\begin{proof}
Let \(P\) be the polar divisor of \(g\), and put
\(\Omega=f^{-1}(B^\circ)\setminus|P|\).  A proper analytic subset of a
connected complex manifold has real codimension at least two and does not
disconnect it.  First note that \(f^{-1}(B^\circ)\) is connected.  Indeed,
each of its connected components is open, and its image is open because
\(f\) is a submersion and closed because the base change
\(f^{-1}(B^\circ)\to B^\circ\) is proper.  Since \(B^\circ\) is connected,
every component would meet every fibre; connectedness of the fibres
therefore permits only one component.  Removing \(|P|\) now leaves
\(\Omega\) connected.  Only finitely many fibres are
components of \(|P|\).  For every other \(b\), the complement
\(F_b\cap\Omega\) of a proper analytic subset in the connected manifold
\(F_b\) is connected.

On \(\Omega\), define the analytic set
\[
 Z=\{dg\wedge f^*(dt)=0\},
\]
where \(t\) is any local coordinate on \(B\).  If \(g|_{F_b}\) is
constant, then \(F_b\cap\Omega\subset Z\).  If \(Z\ne\Omega\), then, because
\(\Omega\) is second countable, \(Z\) has at most countably many
irreducible components.  Each irreducible two-dimensional fibre contained
in \(Z\) is one of those components, and distinct fibres give distinct
components.  This contradicts the uncountable hypothesis.  Therefore
\(Z=\Omega\).

It follows that \(dg\) annihilates \(\ker df\).  Connectedness of
\(F_b\cap\Omega\), together with local holomorphic sections of the
submersion, gives \(g=\varphi\circ f\) on \(\Omega\), for a holomorphic
\(\varphi\) away from finitely many points of \(B\).

Fix such a point \(t_0\).  Choose a general holomorphic curve germ
\(\kappa:(\Delta,0)\to(X,x)\), not contained in \(|P|\), with
\[
 f\circ\kappa(s)=t_0+s^ku(s),\qquad u(0)\ne0.
\]
If \(g\circ\kappa\) has pole order at most \(M\), then
\(|g\circ\kappa(s)|\le C|s|^{-M}\).  After shrinking the disc, boundedness
above and below of \(|u|\) gives constants \(c,C_1>0\) such that
\[
 c|t-t_0|^{1/k}\le |s|\le C_1|t-t_0|^{1/k}
\]
on each local inverse branch.  Therefore
\[
 |\varphi(t)|\le C_2|t-t_0|^{-M/k}
\]
away from the image of a discrete exceptional set, and then everywhere on
the punctured disc by holomorphy.  For any integer \(N>M/k\), the function
\((t-t_0)^N\varphi(t)\) is bounded and extends holomorphically.  Thus
\(\varphi\) is meromorphic at \(t_0\).  Repeating at the finitely many
exceptional points produces \(\varphi\in\mathcal M(B)\).  Equality on the
dense open set \(\Omega\) gives equality of meromorphic functions on \(X\),
and surjectivity of \(f\) gives uniqueness.
\end{proof}

\begin{theorem}[algebraic dimension and algebraic reduction]
\label{thm:algebraic-dimension-audit}
\Verified\ For the global holomorphic family of
\cref{thm:candidate-global-geometry},
\[
 \mathcal M(X)=f^*\mathcal M(\mathbb P^1)=f^*\mathbb C(t),
 \qquad a(X)=1,
\tag{3.9}
\]
and \(f\) is the algebraic reduction.
\end{theorem}

\begin{proof}
By \cref{prop:very-general-torus-audit}, all but countably many smooth
fibres have algebraic dimension zero.  Let \(g\in\mathcal M(X)\).  Only
finitely many fibres are components of its polar divisor, so on uncountably
many of the preceding fibres the restriction of \(g\) is a meromorphic
function on a torus of algebraic dimension zero and is therefore constant.
By \cref{lem:meromorphic-descent-audit}, \(g=f^*\varphi\) for a unique
\(\varphi\in\mathcal M(\mathbb P^1)\).  The reverse inclusion is immediate,
and (3.9) follows.  Since \(f\) is proper, surjective, has connected fibres,
and induces the equality of meromorphic function fields, it is the algebraic
reduction.
\end{proof}

\paragraph{Disposition.}
The earlier label ``algebraic reduction open'' was too broad.  The direct
proof above validates the assertion independently of Section~7 and of the
CDP comparison.  Its provenance is the candidate's Lemmas~9.2--9.6 and
Proposition~9.7, reconstructed here with an explicit Appell--Humbert matrix
bridge, a proof that the smooth locus remains connected, and the omitted
root estimate in the final meromorphic-extension step.  This does not by
itself prove the claimed topology of \(X\); that is a separate calculation.
